% status: ZERO
% Chapter 7 of the second edition. Plan: docs/STRUCTURE.md. Rules: AUTHORING.md.
% Measurements: research/benchmarks/chapter-06-*.md (first edition, scalar CPU),
% research/measurements/2026-09-23-m1-part2.md (Accelerate and Metal).
\chapter{Matrix Multiplication on Real Hardware}\label{ch:matmul-hardware}

\begin{ChapterIntent}
The same multiply--adds can run forty times faster or slower in the same
library on the same GPU, and a thousand times across the ladder from scalar
loops to matrix units, depending on how they are ordered. This chapter follows matrix
multiplication from the scalar loops of Appendix~A onto real hardware: tiles
that live in registers and shared memory, tensor cores that multiply small
matrices per instruction, Hopper's asynchronous warpgroup MMA and Blackwell's
tensor memory. It shows why the shape of a product --- especially the skinny
products of decode --- decides which roof binds and how much of the machine a
kernel can use, and when split-K or a specialized matrix--vector kernel beats a
general GEMM.
\end{ChapterIntent}

\OpeningSection{Same FLOPs, forty times apart}

Multiply two $4{,}096\times4{,}096$ matrices of 16-bit floats on the M1's GPU
and the product --- 137 billion floating-point operations --- takes 62~ms
through PyTorch's MPS (Metal) backend: 2.2~TFLOP/s, about 85\% of the 2.6~TFLOP/s
Apple quotes for the chip (record \texttt{2026-09-23-m1-part2}). Compute the
same product as 4,096 matrix--vector products, one column of the right-hand
matrix at a time, and it takes about 2.5~s (512 of them timed, then scaled):
the same multiply--adds, in the same library, on the same GPU,
\textbf{41 times} slower. Each matrix--vector product reads the whole 32~MiB
left-hand matrix to perform 34 million FLOPs; the single product reads it once
and uses each element 4,096 times.

The ladder below the GPU is longer. This book's first edition timed a scalar
triple loop on the M1's CPU at 2.1~GFLOP/s in the textbook $i,j,k$ order and
19.6 in $i,k,j$ order, at size 256 (Appendix~\ref{app:matmul}). Apple's
Accelerate library runs the $4{,}096$ product on the same CPU at 873~GFLOP/s;
the GPU runs it at 1,604~GFLOP/s in 32-bit floats and 2,213 in 16-bit. A
thousand-fold separates the bottom rung from the top, and none of it comes from
doing less arithmetic. It comes from where the operands are when they are
needed, how many times each byte is used once it is there, and whether the
arithmetic runs on units built for matrices.

One more measurement shows why an engine cannot choose one kernel and stop.
With the same $4{,}096\times4{,}096$ matrix and $N$ right-hand columns, the
library takes 1.77~ms at $N=8$ and 1.02~ms at $N=16$ --- twice the work in 57\%
of the time (73\% in 32-bit floats, which show the same step). Arithmetic
cannot explain that step; a change of kernel between the two shapes can
\todo{verify: capture a
Metal GPU trace to name the kernels selected at $N=8$ and $N=16$}. Decode runs
at exactly these shapes.

\NapkinSection{Napkin math: the work and the bytes of a GEMM}

A product $\mathbf{C}=\mathbf{A}\mathbf{B}$ with
$\mathbf{A}\in\mathbb{R}^{M\times K}$ and $\mathbf{B}\in\mathbb{R}^{K\times N}$
costs $2MNK$ FLOPs and must move at least $(MK+KN+MN)\,b$ bytes. In a decode
step, $\mathbf{A}$ is a weight matrix and $N$ is the batch: the weights
dominate the bytes and the intensity is about $2N/b$ (\Chref{roofline}). That is
the compulsory traffic. What a kernel actually moves depends on how it reuses
data on chip. A kernel that computes an output tile of $T_M\times T_N$ loads
$T_M$ rows of $\mathbf{A}$ and $T_N$ columns of $\mathbf{B}$ for each step
along $K$, and performs $2T_MT_N$ FLOPs per $K$-element for
$(T_M+T_N)\,b$ bytes loaded:
\begin{equation}
I_{\text{tile}} = \frac{2\,T_M\,T_N}{(T_M+T_N)\,b}\quad[\text{FLOP/byte}].
\label{eq:tile-intensity}
\end{equation}
A $128\times128$ tile in BF16 reaches 64~FLOP/byte from slow memory --- short
of an H100's BF16 ridge of about 300 (\Chref{roofline}), which is why fast
kernels also reuse tiles through the L2 cache and share them between
neighbouring blocks. The tile size is bounded by on-chip capacity: registers
for the accumulators, shared memory for the operands, 228~KB per SM on an H100
\cite{nvidia2026hoppertuning}. And at decode, $N$ is the batch: a
$128\times128$ tile with $N=8$ does $8/128$ of the useful work its shape
promises.

\section{From loops to tiles}

Appendix~\ref{app:matmul} builds the product as three nested loops and measures
what the first two optimizations buy on this book's reference machine: changing
the loop order so the innermost loop walks memory contiguously, and blocking the
loops so a working set stays in cache. Both keep the arithmetic identical and
change only which bytes are near the processor when they are needed. Real
kernels continue down the same path --- blocks for the cache, sub-blocks for
shared memory, register tiles for each thread --- until the innermost operation
is a hardware instruction that multiplies a small matrix at once.

\EngineFigure{ch06-hardware}{The same tiled product on two kinds of hardware:
vector lanes of one CPU core working through a tile, and a GPU's cooperating
thread groups each owning part of an output tile. First drawn for
Appendix~\ref{app:matmul}, where the book's scalar kernels stop short of
either.}{fig:ch07-hardware}

\section{The memory hierarchy a tile lives in}

A GPU matrix kernel spans five levels of storage, each much smaller and faster
than the one below: HBM or GDDR (tens to hundreds of gigabytes, a few
terabytes per second), a shared L2 cache (50~MB on an H100), per-SM shared
memory (228~KB on an H100), registers, and --- on Blackwell --- tensor memory
dedicated to matrix accumulators \cite{nvidia2026hoppertuning,nvidia2026cutlassblackwell}.
A well-written kernel stages each operand tile down this hierarchy once and
reuses it as many times as the output tile allows. Appendix~\ref{app:gpu}
gives the numbers for each generation.

\section{Tensor cores: matrix units}

A tensor core executes a small matrix multiply--accumulate, $\mathbf{D}=
\mathbf{A}\mathbf{B}+\mathbf{C}$ on tiles of a few dozen to a few hundred
elements, as a single instruction. It is why a GPU's matrix throughput is an
order of magnitude above its ordinary floating-point throughput, and why
low-precision formats matter: each generation's tensor cores double their rate
at each halving of operand width (\Chref{low-precision-float}). A kernel that
does not feed the tensor cores --- because its operands are the wrong shape,
type or layout --- runs at a fraction of the GPU's advertised peak.

\section{Hopper: asynchronous warpgroup MMA and TMA}

Hopper changed how kernels feed the tensor cores. Its warpgroup MMA
instruction, \texttt{wgmma}, is issued by a warpgroup of 128 threads, reads its
$\mathbf{B}$ operand (and optionally $\mathbf{A}$) directly from shared memory,
keeps the accumulator in registers, and runs asynchronously: a kernel issues
it, does other work, and waits for it with explicit fences. Each instruction
computes a $64\times N\times K$ product with $N$ from 8 to 256
\cite{nvidia2026cutlasswgmma}. Its partner is the Tensor Memory Accelerator,
which copies multi-dimensional tiles between global and shared memory on the
request of a single thread, without spending registers or instruction issue on
address arithmetic \cite{nvidia2026hoppertuning}. Together they let a kernel
dedicate some warps to loading and others to computing, the
\emph{warp specialization} that FlashAttention-3 uses (\Chref{flashattention}).

\FigureIntent{One output tile's journey on Hopper and on Blackwell: TMA copies
operand tiles from HBM into shared memory while producer warps wait on
barriers; consumer warpgroups issue \texttt{wgmma} with the accumulator in
registers (Hopper) or \texttt{tcgen05} with the accumulator in tensor memory
(Blackwell); the epilogue writes the tile back.}

\section{Blackwell: tensor memory and two-SM MMA}

Blackwell's \texttt{tcgen05} instructions are 2--4 times faster than Hopper's
\texttt{wgmma}, depending on data type; they write accumulators into tensor
memory, a new on-chip store, instead of registers; a pair of SMs can execute
one larger MMA together; and block-scaled variants consume FP8, FP6 and FP4
operands with their scale factors directly \cite{nvidia2026cutlassblackwell}.
Freeing registers from accumulators changes kernel design: tiles can be larger,
and the registers go to the non-matrix work --- softmax, dequantization ---
that \Chref{flashattention} found becoming the bottleneck.

\section{Occupancy, waves and tails}

A GEMM's output is divided into tiles, and the tiles are handed to SMs in
waves. If the number of tiles is not a multiple of the number of SMs, the
last wave runs partly empty: a product with 1.1 waves' worth of tiles
leaves most of the GPU idle for the second wave. This \Term{wave quantization}
is invisible in FLOP counts and can cost a large fraction of a kernel's speed
on unlucky shapes. It is also why the same model can run at different
efficiencies on GPUs with different SM counts.

\section{Skinny GEMMs at decode}

At decode, a linear layer multiplies a large weight matrix by $N$ columns,
where $N$ is the number of sequences in the batch: one, eight, sixty-four. A
kernel tuned for square products wastes most of its tile on such a shape,
exactly the tile-quantization effect the M1 showed at 9--32 sequences
(\Chref{roofline}). The other side of the shape is the lever of
\Chref{roofline} seen at kernel level. The opening's $4{,}096\times4{,}096$
product reads its matrix at 36~GB/s for one column; with 16 columns it does 16
times the work in 1.1 times the time, and with 128 columns, 128 times the work
in 2.6 times the time --- until the columns are enough to make the product
compute-bound. Engines therefore carry several kernels for the same
product: matrix--vector kernels for $N$ of one or a few, small-$N$ tilings in
between, and full GEMMs for prefill --- and choose among them by shape at every
step.

\section{Split-K and stream-K}

When $M\times N$ yields too few tiles to fill the GPU --- a skinny product, or
a small layer --- the $K$ dimension can be split instead: several blocks compute
partial sums over different ranges of $K$ for the same output tile, and a final
reduction adds them. This is the same idea as Flash-Decoding's split of the
sequence (\Chref{decode-attention}), and it has the same cost: a reduction pass
and a workspace. Stream-K generalizes it by giving every SM an equal share of
the product's total inner-loop iterations, whatever the tiling, which removes
the wave-quantization loss; its authors reported gains of up to 14 times on some
of 32,824 shapes with a single tile configuration per precision
\cite{osama2023streamk}.

\LedgerSection

Tiling a matrix product for tensor cores, against a straightforward loop:

\begin{ConstraintLedger}
Compute & buys & Access to the matrix units' peak; two to three orders of magnitude over scalar loops on the M1 (measured ladder above). \\
Memory bandwidth & buys & Operand reuse raises the tile's intensity (equation~\eqref{eq:tile-intensity}). \\
Memory capacity & spends & Registers, shared memory and tensor memory for the tiles --- the scarcest capacity on the chip. \\
Latency & risks & At small $N$, a large tile wastes its columns; split-K adds a reduction pass. \\
Correctness & risks & Accumulation precision: FP8 accumulators on Hopper keep about 14 bits (\Chref{low-precision-float}). \\
\end{ConstraintLedger}

\LabSection{tile a GEMM and watch it climb}

\LabIntent{In \texttt{code/labs/ch07-gemm/}, a Rust ladder of matrix
products --- naive, loop-reordered, cache-blocked, register-tiled with
compiler vectorization, and multi-threaded --- each checked against a float64
oracle before timing, plotted in GFLOP/s against the machine's roofline from
the lab of \Chref{roofline}. A second part fixes $M=K=4096$ and varies $N$ from
1 to 512 to reproduce the skinny-GEMM cliff, then adds a split-K variant and
finds the shapes where it wins.}

\HappenedSection{the first edition's loop-order and tiling measurements}

This book's own first edition measured the bottom of the ladder on the M1 and
kept the losses (records in \texttt{research/benchmarks/}, 2026-09-03). Changing
the loop order from $i,j,k$ to $i,k,j$ made a single-threaded scalar product
3.4 to 9.2 times faster at sizes 64 to 256 --- identical arithmetic, contiguous
memory. The blocked kernel's teaching default of $32\times32\times32$ tiles was
not the best tile on that machine: 64 was, at 4.3 times the unblocked loop. The
blocked kernel \emph{lost} at sizes 8 and 16, where its bookkeeping outweighed
its reuse. And with a fixed $512\times512$ weight matrix, eight columns ran at
fewer FLOPs per second through the blocked kernel than one column through the
matrix--vector loop, while sixty-four columns ran fastest of all --- the skinny
product's cliff, measured in scalar code before any GPU was involved. None of
those ratios is a portable constant; all of them are the reason engines choose
kernels by shape.

\FrontierSection{2026-09-24}

Each GPU generation changes the instruction that kernels are built around:
Hopper's asynchronous warpgroup MMA and TMA \cite{nvidia2026cutlasswgmma}, then
Blackwell's tensor memory, paired-SM MMA and native block-scaled low-precision
operands \cite{nvidia2026cutlassblackwell}. Kernel authoring is following with
Python-level tile languages (\Chref{kernel-languages}). \todo{verify: which
stream-K or split-K schedules current CUTLASS and cuBLAS select by default for
decode-shaped products}.

\ExercisesSection

\begin{enumerate}
\item \emph{Derivation.} Using equation~\eqref{eq:tile-intensity}, find the
smallest square tile whose intensity exceeds an H100's BF16 ridge. How much
shared memory would its operands need?
\item \emph{Prediction.} A product has 150 output tiles on a GPU with 132
SMs. Predict its efficiency relative to one with 132 tiles, and propose a
split that recovers it.
\item \emph{Implementation.} Add a split-K path to the lab and find the
smallest $N$ at which the ordinary tiled kernel is faster.
\item \emph{Falsification.} Find a matrix size at which your machine's BLAS is
slower per FLOP than at a smaller size, and explain it.
\end{enumerate}
